\documentclass[10pt]{article}
\usepackage[T1]{fontenc}
\usepackage[utf8]{inputenc}
\usepackage[french]{babel}
\usepackage{amsmath,amssymb,mathrsfs,mathpazo}
\usepackage{geometry}
\geometry{a4paper,margin=2cm}
\usepackage{graphicx,array,tabularx,tabularray,colortbl,multirow,multicol,booktabs,enumitem,pifont,ulem,fancybox,diagbox,cancel}
\usepackage{pstricks,pstricks-add,pst-plot,pst-tree,pst-node,pst-text,pst-func,pst-math,pst-bspline,pst-3dplot,pst-eucl,pst-solides3d}
\let\pscalculate\undefined
\usepackage{pst-bezier}
\usepackage{tikz,pgfplots,tkz-tab}
\usetikzlibrary{arrows,arrows.meta,patterns,calc,positioning,trees,decorations.markings,shapes.geometric,shapes.misc,shapes.arrows,intersections,decorations.pathmorphing,decorations.pathreplacing,backgrounds,shadows,babel}
\pgfplotsset{compat=1.18}
\usepgfplotslibrary{fillbetween}
\usepackage[np]{numprint}
\usepackage{listings,makecell,pas-tableur,twoopt}
\let\barre\undefined
\usepackage[pstricks]{tablvar}
\definecolor{darkgreen}{RGB}{0,100,0}
\providecommand{\red}{\color{red}}
\providecommand{\blue}{\color{blue}}
\providecommand{\green}{\color{green}}
\providecommand{\black}{\color{black}}
\providecommand{\white}{\color{white}}
\providecommand{\gray}{\color{gray}}
\providecommand{\cyan}{\color{cyan}}
\providecommand{\magenta}{\color{magenta}}
\providecommand{\yellow}{\color{yellow}}
\providecommand{\violet}{\color{violet}}
\providecommand{\orange}{\color{orange}}
\providecommand{\euro}{\texteuro}
\providecommand{\EUR}{\texteuro}
\providecommand{\ieme}{\textsuperscript{e}}
\providecommand{\degres}{\ensuremath{^\circ}}
\providecommand{\degree}{\ensuremath{^\circ}}
\providecommand{\vv}[1]{\overrightarrow{#1}}
\providecommand{\Coord}[1]{#1}
\providecommand{\decosix}{}
\providecommand{\starredbullet}{\ensuremath{\bullet}}
\pagestyle{empty}
\setlength{\parindent}{0pt}
\usepackage{adjustbox}\DeclareUnicodeCharacter{2236}{:}\DeclareUnicodeCharacter{2190}{\ensuremath{\leftarrow}}
\providecommand{\R}{\mathbb{R}}
\providecommand{\N}{\mathbb{N}}
\providecommand{\Z}{\mathbb{Z}}
\providecommand{\Q}{\mathbb{Q}}
\providecommand{\D}{\mathbb{D}}
\providecommand{\C}{\mathbb{C}}
\providecommand{\e}{{\mathrm{e}}}
\providecommand{\vect}[1]{\overrightarrow{#1}}
\providecommand{\vectt}[1]{\overrightarrow{\mathrm{#1}}}
\providecommand{\vvt}[1]{\overrightarrow{\mathrm{#1}}}
\providecommand{\np}[2][]{{#2~#1}}
\providecommand{\vv}[1]{{\overrightarrow{#1}}}
\renewcommand{\d}{{\mathrm{d}}}
\providecommand{\degres}{{^\circ}}
\providecommand{\degre}{{^\circ}}
\providecommand{\blue}{\color{blue}}
\providecommand{\red}{\color{red}}
\providecommand{\enskip}{\quad}
\providecommand{\llbracket}{\lbrack\!\lbrack}
\providecommand{\rrbracket}{\rbrack\!\rbrack}
\providecommand{\euro}{{\text{€}}}
\providecommand{\no}{{\text{n°}}}
\providecommand{\interval}[2]{{[#1;#2]}}
\providecommand{\intervFO}[2]{{[#1;#2[}}
\providecommand{\dotfill}{\ldots}

\providecommand{\ds}{\displaystyle}
\providecommand{\textding}[1]{#1}
\providecommand{\Oij}{\(\mathrm{O};\vec{\imath},\vec{\jmath}\)}
\providecommand{\Oijk}{\(\mathrm{O};\vec{\imath},\vec{\jmath},\vec{k}\)}
\providecommand{\Ouv}{\(\mathrm{O};\vec{u},\vec{v}\)}
\providecommand{\intervFF}[2]{\ensuremath{[#1;#2]}}
\renewcommand{\euro}{\ifmmode\text{\texteuro}\else\texteuro\fi}
\DeclareUnicodeCharacter{00B0}{\ifmmode^{\circ}\else\textdegree\fi}
\DeclareUnicodeCharacter{20AC}{\ifmmode\text{\texteuro}\else\texteuro\fi}
% Énoncés : Éducation nationale. Transcription et corrigés : APMEP, Denis Vergès et contributeurs. Découpage et classement : CoopMaths, Sébastien Lozano, Éric Elter et contributeurs. Adaptation de la présentation : Maths Explorer.
% CC BY-SA : https://coopmaths.fr/www/about/
% https://coopmaths.fr/alea/static/bac/2023/tex/bac_2023_09_sujet1_ameriquesud_2.tex
\begin{document}
\noindent\textit{APMEP / CoopMaths. CC BY-SA. Présentation adaptée par Maths Explorer.}\par\medskip
\medskip

\begin{enumerate}
\item Entre 1998 et 2020, en France \np{18221965} accouchements ont été recensés, parmi lesquels
\np{293898} ont donné naissance à des jumeaux et \np{4921} ont donné naissance à au moins trois enfants.
	\begin{enumerate}
		\item Avec une précision de 0,1\,\% calculer parmi tous les accouchements recensés, le pourcentage d'accouchements donnant naissance à des jumeaux sur la période 1998-2020.
		\item Vérifier que le pourcentage d'accouchements qui ont donné naissance à au moins trois enfants est inférieur à 0,1\,\%.
\end{enumerate}
\end{enumerate}

On considère alors que ce pourcentage est négligeable.

On appelle accouchement ordinaire, un accouchement donnant naissance à un seul enfant.

On appelle accouchement double, un accouchement donnant naissance à exactement deux enfants.

On considère dans la suite de l'exercice qu'un accouchement est soit ordinaire, soit double.

La probabilité d'un accouchement ordinaire est égale à $0,984$ et celle d'un accouchement double est alors égale à $0,016$.

Les probabilités calculées dans la suite seront arrondies au millième.

\begin{enumerate}[resume]
\item On admet qu'un jour donné dans une maternité, on réalise $n$ accouchements.

On considère que ces $n$ accouchements sont indépendants les uns des autres.

On note $X$ la variable aléatoire qui donne le nombre d'accouchements doubles pratiqués ce
jour.
	\begin{enumerate}
		\item Dans le cas où $n =  20$, préciser la loi de probabilité suivie par la variable aléatoire $X$ et calculer la probabilité qu'on réalise exactement un accouchement double.
		\item Par la méthode de votre choix que vous expliciterez, déterminer la plus petite valeur de $n$ telle que $P(X \geqslant 1) \geqslant  0,99$.

Interpréter le résultat dans le contexte de l'exercice.
	\end{enumerate}
\item Dans cette maternité, parmi les naissances double, on estime qu'il y a 30\,\% de jumeaux monozygotes (appelés \og  vrais jumeaux \fg{}  qui sont obligatoirement de même sexe : deux garçons ou deux filles) et donc 70\,\% de jumeaux dizygotes (appelés \og  faux jumeaux \fg, qui peuvent être de sexes différents: deux garçons, deux filles ou un garçon et une fille).

Dans le cas de naissances doubles, on admet que, comme pour les naissances ordinaires, la probabilité d'être une fille à la naissance est égale à $0,49$ et que celle d'être un garçon à la naissance est égale à $0,51$.

Dans le cas d'une naissance double de jumeaux dizygotes, on admet aussi que le sexe du
second nouveau-né des jumeaux est indépendant du sexe du premier nouveau-né.

On choisit au hasard un accouchement double réalisé dans cette maternité et on considère les évènements suivants :

\begin{itemize}
\item[$\bullet~~$]$M$ : \og  les jumeaux sont monozygotes \fg{}  ;
\item[$\bullet~~$]$F_1$ : \og le premier nouveau-né est une fille \fg{} ;
\item[$\bullet~~$]$F_2$ : \og le second nouveau-né est une fille \fg.
\end{itemize}

\emph{On notera $P(A)$ la probabilité de l'évènement $A$ et $\overline{A}$ l'évènement contraire de $A$.}

\begin{center}
\pstree[treemode=R,nodesepA=0pt,nodesepB=2.5pt,treesep = 1cm,levelsep=2.5cm]{\TR{}}
{\pstree{\TR{$M$~~}}
	{\pstree{\TR{$F_1$}}
		{\TR{$F_2$}	
		}
	\pstree{\TR{$\overline{F_1}$~}}
		{\TR{$\overline{F_2}$}	
		}
	}
\pstree{\TR{$\overline{M}$~~}}
	{\pstree{\TR{$F_{1}$~}}
		{\TR{$F_2$}
		\TR{$\overline{F_2}$}
		}
	\pstree{\TR{$\overline{F_1}$~}}
		{\TR{$F_2$}
		\TR{$\overline{F_2}$}
		}
	}
}
\end{center}
\smallskip

	\begin{enumerate}
		\item Recopier puis compléter l'arbre pondéré ci-dessus.
		\item Montrer que la probabilité que les deux nouveaux-nés soient des filles est \np{0,31507}.
		\item Les deux nouveaux-nés sont des jumelles. Calculer la probabilité qu'elles soient monozygotes.
	\end{enumerate}
\end{enumerate}



\bigskip

\end{document}
